\documentclass[runningheads]{llncs}
%
\usepackage[T1]{fontenc}
\usepackage{graphicx}
\usepackage{amsmath,amssymb}
\usepackage{booktabs}
\usepackage{multirow}
\usepackage{array}
%
\emergencystretch=3em
\begin{document}
%
\title{GIS-Based Smart Village Infrastructure Planning Using DINOv2 Segmentation, Graph-Based Connectivity Analysis, and A$^{*}$ Route Planning}
%
\titlerunning{GIS-Based Smart Village Infrastructure Planning}
%
\author{Ramakrishna Kolikipogu\and Pendyala Archana\and Palamakula Swathi}
%
\authorrunning{R. Kolikipogu et al.}
%
\institute{Department of Information Technology, Chaitanya Bharathi Institute of Technology (A), Hyderabad, India\\
\email{krkrishna.cse@gmail.com}\\
\email{pendyalaarchanareddy98@gmail.com}\\
\email{swathip\_it@cbit.ac.in}}
%
\maketitle
%
\begin{abstract}
Rural infrastructure planning matters for living standards, but field-survey approaches are slow and hard to scale. This paper presents StrongDINOv2Seg, coupling a DINOv2 vision-transformer encoder with an FPN, Atrous Spatial Pyramid Pooling (ASPP), and a residual decoder to segment LoveDA imagery into its seven semantic classes (Background, Building, Road, Water, Barren, Forest, Agriculture; the no-data index is ignored). Trained for 10 epochs on combined Rural/Urban LoveDA, the model reaches a best validation mIoU of 53.80\% and mDice of 68.91\% at epoch 4 (no test-set evaluation was performed). An OpenCV/NetworkX/SciPy stage then performs building detection, road skeletonization, graph construction, KD-tree connectivity, and A$^{*}$ route planning over a land-cover-weighted cost surface, generating raster-based computational connectivity proposals for isolated buildings. On a representative tile with 14 buildings, connectivity rises from 28.57\% to 100\% via 10 proposed segments (2,029.94 px total, path efficiency 0.86); over ten further tiles (54 buildings), connectivity rises from 40.74\% to 96.30\% via 30 connections (efficiency 0.931). All figures are validation-split results with no held-out test set; planning covers only the selected LoveDA tiles; and route outputs are decision-support proposals, not field-validated and not establishing engineering, terrain, environmental, cost, legal, or social feasibility.
\keywords{DINOv2 \and Semantic segmentation \and ASPP \and FPN \and LoveDA dataset \and Smart village \and Rural infrastructure planning \and Road extraction \and Building detection \and Connectivity analysis \and KD-tree \and A$^{*}$ path planning \and OpenCV}
\end{abstract}
%
\section{Introduction}
Characterising rural infrastructure -- roads, buildings, water, agricultural land -- matters for living standards~\cite{hinton1996}. Poor road connectivity is a recognised barrier to healthcare access~\cite{verma2020}, and infrastructural gap analysis has been used to prioritise rural investment~\cite{kapoor2021}. GIS-based imagery analysis supports this, but manual survey is slow and hard to scale; benchmarks such as SpaceNet~\cite{vanetten2018} and DeepGlobe~\cite{demir2018} have driven progress in automated extraction, motivating its combination with planning.

This paper describes StrongDINOv2Seg: a DINOv2-based segmentation model paired with a downstream pipeline -- road extraction, skeletonization, graph construction, KD-tree connectivity, and A$^{*}$ route planning -- that turns segmentation into raster-based computational connectivity proposals. It is evaluated on LoveDA segmentation quality, a representative planning tile, and a ten-tile sample, reporting class-wise metrics, an A$^{*}$ cost formulation, and a radius-sensitivity analysis.

\section{Related Work}
\textbf{Segmentation.} SegFormer~\cite{xie2021} and UNetFormer~\cite{wang2022unet} (on LoveDA) advanced transformer segmentation but need large labelled sets; DINOv2~\cite{oquab2023} is self-supervised but not built for dense satellite classification, addressed here via fine-tuning with an ASPP/FPN decoder. \textbf{Road/building extraction.} D-LinkNet~\cite{zhou2018} won the DeepGlobe road-extraction challenge; attention, atrous/strip, fusion, and dual-decoder extractors followed~\cite{bai2025,ma2026,zhang2024,qu2025}, with several targeting road topology~\cite{kamada2021,mattyus2017,xu2023,liu2017}. Ayala et al.~\cite{ayala2021} jointly detected buildings and roads from Sentinel imagery, mirroring our two-class task; VHR building/settlement classification~\cite{dimassi2024,kakooei2026}, a modified ViT for rural land cover~\cite{wassan2025}, and MAP-Net~\cite{zhu2021} for small buildings are also relevant. \textbf{GIS planning and routing.} OSMnx~\cite{boeing2017} popularised the NetworkX street-graph used in our connectivity stage; Verma and Dash~\cite{verma2020} modelled rural healthcare accessibility via GIS network analysis, and Kapoor et al.~\cite{kapoor2021} ranked Telangana village infrastructure gaps, both relevant to prioritising isolated buildings. Wang et al.~\cite{wang2022astar} proposed weighted, cost-aware A$^{*}$ routing, related to our traversal-cost surface. These works generally treat segmentation, extraction, or GIS accessibility in isolation.

\textbf{Problem and contributions.} A pipeline that only classifies land cover does not, by itself, tell a planner where new access roads are needed. Planning requires a rural land-cover segmenter, a discrete inventory of buildings and roads, and a rule for deciding which buildings need access and where; prior systems address these separately. This paper integrates them as an engineering decision-support demonstration whose outputs are raster-based computational connectivity proposals, not field-validated engineering plans. The main contributions are: (i) integrating a DINOv2-based segmentation model with an explicit cost-weighted A$^{*}$ route-planning stage; (ii) class-wise and aggregate segmentation metrics; (iii) an explicit A$^{*}$ traversal-cost surface; (iv) evaluation on a representative tile, a radius-sensitivity sweep, and a pooled ten-tile sample; and (v) explicit limitations. No component is individually novel; the contribution is their integration.

\section{Proposed Methodology}
\subsection{Overall System Architecture}
The pipeline has two stages: StrongDINOv2Seg classifies each pixel into one of eight indices per the official LoveDA encoding (0 = no-data, ignored; 1--7 = Background, Building, Road, Water, Barren, Forest, Agriculture); a second stage then uses classical vision and graph algorithms to detect buildings, skeletonize roads, mark existing connections, and plan segments for isolated ones (Sect.~\ref{sec:setup}).

\subsection{DINOv2 Encoder, FPN, ASPP, and Decoder}
A pretrained \texttt{facebook/dinov2-base} Vision Transformer with patch size 14 serves as the fine-tuned encoder. Inputs are resized to $560\times560$ and embedded into $40\times40$ patch tokens. Hidden states are taken from layers 3, 6, 9, and 12 (768-dimensional each):
\begin{equation}
z'_l=\mathrm{MSA}(\mathrm{LN}(z_{l-1}))+z_{l-1},\qquad z_l=\mathrm{MLP}(\mathrm{LN}(z'_l))+z'_l
\end{equation}
where MSA is multi-head self-attention, LN is layer normalisation, and MLP is the feed-forward block. Each $40\times40\times768$ feature map is projected to 256 channels ($P_l=\mathrm{Conv}_{1\times1}(f_l)$), concatenated, and fused:
\begin{equation}
F_{\mathrm{fuse}}=\mathrm{Conv}_{1\times1}\big(\mathrm{Concat}(P_3,P_6,P_9,P_{12})\big)
\end{equation}
ASPP is applied with dilation rates 6, 12, 18, a $1\times1$ branch, and a global-pooling branch, giving a 256-channel output: $\mathrm{ASPP}(x)=\mathrm{Proj}(\mathrm{Concat}(B_1,\dots,B_5))$. A residual decoder ($R(x)=F(x)+x$) restores spatial resolution through four $2\times$ upsampling stages ($256\to256\to128\to64\to32$ channels), producing a $640\times640$ map resized to $560\times560$, followed by a final $1\times1$ convolution to eight classes.

\subsection{Loss Function}
Training uses a composite loss:
\begin{equation}
L_{\mathrm{total}}=0.5\,L_{\mathrm{CE}}+0.4\,L_{\mathrm{Dice}}+0.1\,L_{\mathrm{boundary}},\qquad L_{\mathrm{Dice}}=1-\frac{2|P\cap G|+\epsilon}{|P|+|G|+\epsilon}
\end{equation}
where $L_{\mathrm{CE}}$ is pixel-wise cross-entropy and $L_{\mathrm{boundary}}$ is a Sobel-edge-weighted boundary term; both CE and Dice ignore the no-data index (0). An auxiliary deep-supervision loss weight is defined but inactive, since the decoder exposes no auxiliary heads.

\section{System Architecture: Connectivity Analysis and Route Planning}
\subsection{Building/Road Extraction and Connectivity Graph}
Buildings are detected via 8-connected component analysis of the Building mask (OpenCV \texttt{connectedComponentsWithStats}), discarding components under 20 px. Road pixels are cleaned by $3\times3$ closing then opening and skeletonized (scikit-image) to a one-pixel centreline, converted to an undirected NetworkX graph -- the OSMnx~\cite{boeing2017} street-graph representation -- with 8-connected edges weighted 1.0 (orthogonal) or $\sqrt2$ (diagonal). A building counts as connected if the road mask, dilated by radius $r$ (default 5 px), overlaps its expanded bounding box; for isolated buildings a KD-tree over skeleton pixels returns the nearest road pixel, or none if the skeleton is empty.

\subsection{Cost-Weighted A$^{*}$ Route Planning}
For every isolated building with a valid KD-tree target, A$^{*}$ search plans a route over an 8-connected grid, following the weighted-cost formulation adopted in GIS route-planning studies such as~\cite{wang2022astar}. Let $I(p)$ denote the predicted class index at pixel $p$:
\begin{equation}
\mathrm{cost}(p)=\begin{cases}
\infty & I(p)\in\{0,4\}\ \text{(no-data, water)}\\
1 & I(p)=3\ \text{(road)}\\
2 & I(p)=7\ \text{(agriculture)}\\
3 & I(p)=1\ \text{(background)}\\
4 & I(p)=5\ \text{(barren)}\\
5 & I(p)=2\ \text{(building)}\\
8 & I(p)=6\ \text{(forest)}
\end{cases}
\end{equation}
The search minimises $F(n)=g(n)+h(n)$ with $g$ the accumulated cost and $h$ the Euclidean heuristic; since all finite costs are $\ge1$, $h$ never overestimates, so a returned route is minimum-cost under this raster model, not a globally optimal real-world design. If A$^{*}$ fails, a straight Bresenham fallback is used; it ignores the cost surface and may cross infinite-cost pixels (e.g., water), so fallback routes are topological placeholders requiring manual re-routing, not traversable proposals. Lengths are reported in pixels on the $560\times560$ grid. LoveDA imagery has a native ground sampling distance of about $0.3$\,m at $1024\times1024$~\cite{wang2021loveda}, so after resizing the effective GSD is roughly 0.55\,m/px, e.g.\ the 202.99 px mean proposed length is approximately 110\,m; these conversions are approximate and unused in any reported metric. Routes remain pixel-space decision-support proposals, not field validated, and establish no engineering, terrain, environmental, cost, legal, or social feasibility.

\begin{table}[t]
\caption{Training configuration.}\label{tab:train}
\centering
\begin{tabular}{llll}
\toprule
Parameter & Value & Parameter & Value\\
\midrule
Encoder & DINOv2-base (p14) & Optimizer & AdamW\\
Resolution/classes & $560^2$ / 8 & Weight decay & $1\times10^{-4}$\\
Batch/epochs & 2 / 10 & Grad. clip/AMP & 1.0 / off\\
Enc./dec. LR & $5\times10^{-6}$ / $5\times10^{-5}$ & Loss (CE/Dice/bnd.) & .5/.4/.1\\
LR schedule & Cosine, $\eta_{\min}=10^{-7}$ & Parameters & 93.8M\\
\bottomrule
\end{tabular}
\end{table}

